\documentclass[10pt,a4paper]{amsart}
\usepackage[margin=19mm]{geometry}
\usepackage{amsmath,amssymb,mathtools}
\usepackage[scr=rsfs,cal=euler]{mathalfa}
\usepackage{xfrac}
\usepackage[unicode,hidelinks,bookmarks=false]{hyperref}
\newcommand{\ie}{i.e.\ }
\newcommand{\cf}{cf.\ }
\newcommand{\resp}{respectively\ }
\newcommand{\Subsection}{\subsection}
\providecommand{\nobreakdash}{\nobreak\hbox{-}}
\setlength{\emergencystretch}{2em}
\multlinegap0pt
\usepackage{bm}
\usepackage{bbm}
\usepackage{dsfont}
\usepackage{xspace}
\usepackage{cancel}
\usepackage{mathtools}
\usepackage{amsthm}
\makeatletter
\@ifundefined{theorem}{\newtheorem{theorem}{Theorem}}{}
\makeatother
\title[An annihilation-number Caro-Wei bound: a TxGraffiti conjectu]{An annihilation-number Caro-Wei bound: a TxGraffiti conjecture and an independence-number bracket}
\author{Chakshu Gupta}
\date{Source: arXiv:2606.29553v1 (CC BY 4.0)}
\begin{document}
\maketitle

\noindent\emph{Source and licence.} An annihilation-number Caro-Wei bound: a TxGraffiti conjecture and an independence-number bracket. Version arXiv:2606.29553v1 (CC BY 4.0), distributed under the Creative Commons Attribution 4.0 International licence, as recorded in the arXiv licence metadata for this exact version and shown on the abstract page https://arxiv.org/abs/2606.29553v1. Full rights notice: licenses/arxiv-2606.29553-CC-BY-4.0.txt.\par\smallskip

\noindent\textit{Editorial scope.} Counted unit: the Theorem whose source label is \texttt{thm:vehicle} in the article (source0.tex lines 111--116, source label \texttt{thm:vehicle}), delivered together with the article's own complete original proof of it (source0.tex lines 118--169, one \texttt{proof} environment of 52 lines). No mathematics is written, repaired or re-derived: statement and proof are the article's own text line for line. Required context retained uncounted: none. Every definition, display and label of those lines is retained unchanged. Omitted from this package and not redistributed: 1--110, 117--117, 170--293. Nothing inside a retained excerpt is omitted or altered: every retained body is delivered exactly as the article's lines are written, so no omission receipt of any kind accompanies this package. Citations: the retained excerpts carry no citation command at all, so no bibliography entry of the article is transcribed and no reference is redistributed. Reference adaptation: none was needed and none was made; every reference inside the delivered unit resolves inside it, so no unresolved reference or citation is produced. Printed slips kept literal: no printed word, symbol, hypothesis or number of the counted statement or of its proof was altered. Layout and class: the article's own class is \texttt{article} with packages including geometry, amsmath, amssymb, amsthm, xcolor, hyperref, url; the wrapper is the lane's pinned \texttt{amsart} editorial wrapper at 10pt on A4, which restates the theorem-like environments the retained text needs and adds the article's own macro definitions that the retained text needs (none), with extra wrapper packages (bm, bbm, dsfont, xspace, cancel, mathtools). The delivered numbering is the unit's own. Assets: no retained span contains a figure environment, a graphics-inclusion command, a PDF-page inclusion, a table or a TikZ picture; no image file is copied into this package, the package declares no asset dependency and no image file is redistributed with it. Licence: version arXiv:2606.29553v1 of this article is distributed under the Creative Commons Attribution 4.0 International licence, as recorded in the arXiv licence metadata for this exact version and shown on the abstract page https://arxiv.org/abs/2606.29553v1. A full rights notice is delivered at licenses/arxiv-2606.29553-CC-BY-4.0.txt. Characters: every retained excerpt is pure ASCII, so the delivered engine, XeLaTeX with its default Latin Modern text font, reports no missing character.
\begin{theorem}\label{thm:vehicle}
Every graphic degree sequence with maximum degree $\Delta \ge 1$ satisfies
\[
  a \ \le\ \frac{\Delta+1}{2}\, W.
\]
\end{theorem}

\begin{proof}
Let $H = \{1, \ldots, a\}$ index the $a$ smallest degrees and
$T = \{a+1, \ldots, n\}$ the remaining $n-a$; call these the head and the tail.

\medskip\noindent\emph{Head degree-sum bound.}
The annihilation number $a$ is the largest $j$ with $d_1 + \cdots + d_j \le m$, so
the head satisfies $\sum_{i \in H} d_i \le m$. Each edge contributes $2$ to the
degree sum, hence $\sum_{i=1}^{n} d_i = 2m$; as the head and tail partition the
indices,
\[
  \sum_{i \in T} d_i \ =\ 2m - \sum_{i \in H} d_i \ \ge\ 2m - m \ =\ m
    \ \ge\ \sum_{i \in H} d_i .
\]
Every degree is at most $\Delta = d_n$, so the $n-a$ tail terms sum to at most
$\Delta(n-a)$. Chaining this with the previous line,
\begin{equation}\label{eq:headsum}
  \sum_{i \in H} d_i \ \le\ \sum_{i \in T} d_i \ \le\ \Delta\,(n-a).
\end{equation}

\medskip\noindent\emph{Pointwise inequality.}
For every integer $k$ with $0 \le k \le \Delta$,
\begin{equation}\label{eq:pointwise}
  \frac{1}{k+1} + \frac{k}{\Delta(\Delta+1)} \ \ge\ \frac{2}{\Delta+1}.
\end{equation}
Multiplying through by the positive number $\Delta(\Delta+1)(k+1)$ and collecting
terms,~\eqref{eq:pointwise} is equivalent to
\[
  \Delta(\Delta+1) + k(k+1) - 2\Delta(k+1) \ \ge\ 0,
\]
whose left-hand side equals $(\Delta-k)(\Delta-k-1)$. Since $0 \le k \le \Delta$,
the integer $\Delta - k$ is nonnegative: if $\Delta - k = 0$, the product is
$0 \cdot (-1) = 0$; if $\Delta - k \ge 1$, then $\Delta - k - 1 \ge 0$ and the
product is the product of two nonnegative integers. In either case it is
nonnegative, with equality precisely when one factor vanishes, that is, at
$k = \Delta$ and $k = \Delta - 1$.

\medskip\noindent\emph{Summation.}
Split $W$ over the head and tail. On the tail every degree is at most $\Delta$, so
$\tfrac{1}{d_i+1} \ge \tfrac{1}{\Delta+1}$, and the $n-a$ tail terms contribute at
least $(n-a)/(\Delta+1)$. On the head apply~\eqref{eq:pointwise} with $k = d_i$,
legitimate as each $d_i \le \Delta$:
\[
  W \ \ge\ \sum_{i \in H}\!\left(\frac{2}{\Delta+1} - \frac{d_i}{\Delta(\Delta+1)}\right)
        + \frac{n-a}{\Delta+1}
    \ =\ \frac{2a}{\Delta+1}
        - \frac{\sum_{i \in H} d_i}{\Delta(\Delta+1)}
        + \frac{n-a}{\Delta+1}.
\]
By~\eqref{eq:headsum}, $\sum_{i \in H} d_i \le \Delta(n-a)$, so the middle term is
at least $-(n-a)/(\Delta+1)$; it cancels the tail contribution, leaving
$W \ge 2a/(\Delta+1)$. Rearranging gives $a \le \tfrac{\Delta+1}{2}W$.
\end{proof}

\end{document}
